\uuid{HfzV}
\titre{Limite}
\auteur{bodin ; hanani ; mzari}
\organisation{exo7}
\licence{CC BY-NC-SA 4.0}
\datecreate{2020-01-20}
\niveau{L1}
\module{Analyse}
\chapitre{Continuité, limite et étude de fonctions réelles}
\sousChapitre{Limite de fonctions}
\difficulte{1}
\theme{Limites des fonctions réelles, limite}
\exo7qcmid{lille-1.limites.004}

\contenu{
\begin{question}
\qclassification{123.03}
\qtags{motcle=limite}
Soit $f(x)= \frac{\sqrt {x+1}-\sqrt {2x}}{x-1}$. Quelles sont les assertions vraies ?
\begin{answers}
    \bad{$\lim_{x\to +\infty} f(x)=+\infty$}

    \good{$\lim_{x\to +\infty} f(x)=0$}

    \good{$\lim_{x\to 1} f(x)=-\frac{1}{2\sqrt 2}$}

    \bad{$f$ n'admet pas de limite en $1$.}
\end{answers}
\begin{explanations}
On multiplie le numérateur et le dénominateur de $f$ par l'expression conjuguée de $\sqrt {x+1}-\sqrt {2x}$, c'est-à-dire par  $\sqrt {x+1}+\sqrt {2x}$. On obtient :
$f(x)= -\frac{1}{\sqrt {x+1}+\sqrt {2x}}$.
\end{explanations}
\end{question}
}
